\begin{bpnew}{propborder}{Definition: $M$-Tail of a Sequence}
\begin{definition}[$M$-Tail of a Sequence]
\label{def:m-tail}

Let $(x_n) \subseteq \mathbb{R}$ and let $M \in \mathbb{N}$. The
\textbf{$M$-tail} of $(x_n)$ is the sequence
\[
    (x_n)_{n \ge M} \;:=\; \bigl(x_M,\, x_{M+1},\, x_{M+2},\, \ldots\bigr).
\]
\end{definition}
\end{bpnew}

\begin{remark*}[Standard quantified statement]
\[
\forall (x_n) \subseteq \mathbb{R} \;\forall M \in \mathbb{N}
\;\Bigl(
    (x_n)_{n \ge M} = \{x_n : n \ge M\}
\Bigr).
\]
\end{remark*}

\begin{remark*}[Interpretation]
The $M$-tail discards the first $M - 1$ terms and retains everything
from index $M$ onward. It is the formal object underlying the phrase
\emph{all but finitely many terms}: to say that all but finitely many
terms of $(x_n)$ satisfy a property $P$ is precisely to say that some
$M$-tail of $(x_n)$ satisfies $P$ termwise. The index $M$ is an offset
into the sequence, not a bound on the terms themselves; two sequences
with identical $M$-tails share all asymptotic properties, differing
only in finitely many initial values which convergence arguments
discard entirely.
\end{remark*}



\begin{bpnew}{thmborder}{Theorem: Convergence of a Tail}
\begin{theorem}[Convergence of a Tail]
\label{thm:convergence-of-tail}

Let $(x_n) \subseteq \mathbb{R}$ and let $m \in \mathbb{N}$. The
$m$-tail $(x_{m+n})_{n \in \mathbb{N}}$ converges if and only if
$(x_n)$ converges. Moreover, in this case,
\[
    \lim_{n \to \infty} x_{m+n} \;=\; \lim_{n \to \infty} x_n.
\]

\smallskip
\noindent
\hyperref[prf:convergence-of-tail]{\textit{Go to proof.}}
\end{theorem}
\end{bpnew}

\begin{remark*}[Standard quantified statement]
\[
\begin{aligned}
&\forall (x_n) \subseteq \mathbb{R} \;\forall m \in \mathbb{N}\\
&\quad\Bigl(
    \bigl(x_{m+n}\bigr)_{n \in \mathbb{N}} \text{ converges}
    \;\Longleftrightarrow\;
    (x_n) \text{ converges}
\Bigr)\\
&\quad\land\;\Bigl(
    \lim_{n\to\infty} x_{m+n} = \lim_{n\to\infty} x_n
\Bigr).
\end{aligned}
\]
\end{remark*}

\begin{remark*}[Contrapositive quantified statement]
\[
\begin{aligned}
&\forall (x_n) \subseteq \mathbb{R} \;\forall m \in \mathbb{N}\\
&\quad\Bigl(
    \bigl(x_{m+n}\bigr)_{n \in \mathbb{N}} \text{ diverges}
    \;\Longleftrightarrow\;
    (x_n) \text{ diverges}
\Bigr).
\end{aligned}
\]
\end{remark*}

\begin{remark*}[Interpretation]
Convergence is a \emph{tail property}: it depends only on the
eventual behavior of the sequence, not on any finite initial
segment. Discarding the first $m$ terms neither creates nor
destroys convergence, and when the limit exists it is preserved
exactly. This theorem formalizes the informal principle that
finitely many terms are irrelevant to asymptotic behavior. It is
the rigorous justification for the common proof move of assuming
without loss of generality that $n$ is large enough for some
auxiliary condition to hold — one simply passes to a suitable
$m$-tail, applies the argument there, and concludes for the full
sequence by this theorem. The contrapositive is equally useful:
to show $(x_n)$ diverges it suffices to show that some $m$-tail
diverges.
\end{remark*}









\begin{bpnew}{thmborder}{Theorem: Convergence by Domination}
\begin{theorem}[Convergence by Domination]
\label{thm:convergence-by-domination}

Let $(x_n) \subseteq \mathbb{R}$ and let $L \in \mathbb{R}$. Let
$(a_n)$ be a sequence of positive real numbers with
$\displaystyle\lim_{n \to \infty} a_n = 0$. If there exist a constant
$c > 0$ and an index $m \in \mathbb{N}$ such that
\[
    |x_n - L| \;\le\; c\, a_n \qquad \forall n \ge m,
\]
then $\displaystyle\lim_{n \to \infty} x_n = L$.

\smallskip
\noindent
\hyperref[prf:convergence-by-domination]{\textit{Go to proof.}}
\end{theorem}
\end{bpnew}

\begin{remark*}[Standard quantified statement]
\[
\begin{aligned}
&\forall (x_n) \subseteq \mathbb{R} \;\forall L \in \mathbb{R}
 \;\forall (a_n) \subseteq (0, \infty) \;\forall c > 0
 \;\forall m \in \mathbb{N}\\
&\quad\Bigl(
    \Bigl(
        \lim_{n \to \infty} a_n = 0
        \;\land\;
        \forall n \ge m \;\bigl(|x_n - L| \le c\, a_n\bigr)
    \Bigr)
    \;\Longrightarrow\;
    \lim_{n \to \infty} x_n = L
\Bigr).
\end{aligned}
\]
\end{remark*}

\begin{remark*}[Contrapositive quantified statement]
\[
\begin{aligned}
&\forall (x_n) \subseteq \mathbb{R} \;\forall L \in \mathbb{R}
 \;\forall (a_n) \subseteq (0,\infty) \;\forall c > 0
 \;\forall m \in \mathbb{N}\\
&\quad\Bigl(
    \lim_{n \to \infty} x_n \neq L
    \;\Longrightarrow\;
    \lim_{n \to \infty} a_n \neq 0
    \;\lor\;
    \exists n \ge m \;\bigl(|x_n - L| > c\, a_n\bigr)
\Bigr).
\end{aligned}
\]
\end{remark*}

\begin{remark*}[Interpretation]
The theorem is a \emph{domination principle}: it is not necessary to
verify the $\varepsilon$-$N$ condition directly for $(x_n)$. It
suffices to find a null sequence $(a_n)$ — one converging to zero —
that dominates the error $|x_n - L|$ on some tail, up to a fixed
scalar $c$. The scalar $c$ is irrelevant to the conclusion: since
$a_n \to 0$, the product $c\, a_n \to 0$ regardless of the size of
$c$, and any $\varepsilon > 0$ is eventually beaten by $c\, a_n$.
This theorem is the rigorous engine behind the common proof pattern
of bounding $|x_n - L|$ by a simpler expression and observing that
the simpler expression tends to zero. The canonical instance is
$a_n = 1/n$, where domination by $c/n$ is sufficient to conclude
convergence to $L$.
\end{remark*}

